%%%%%%%%%%%%%%%%%%%%%%%%%%%%%%%%%%%%%%%%%%%%%%%%%%%%%%%
% Arquivo para entrada de dados para a parte pré textual
%%%%%%%%%%%%%%%%%%%%%%%%%%%%%%%%%%%%%%%%%%%%%%%%%%%%%%%
%
% Basta digitar as informações indicidas, no formato
% apresentado.
%
%%%%%%%
% Os dados solicitados são, na ordem:
%
% tipo do trabalho
% componentes do trabalho
% título do trabalho
% nome do autor
% local
% data (ano com 4 dígitos)
% orientador(a)
% coorientador(a)(as)(es)
% arquivo com dados bibliográficos
% instituição
% setor
% programa de pós gradução
% curso
% preambulo
% data defesa
% CDU
% errata
% assinaturas - termo de aprovação
% resumos & palavras chave
% agradecimentos
% dedicatoria
% epígrafe


% Informações de dados para CAPA e FOLHA DE ROSTO
%-----------------------------------------------------------------------------
\tipotrabalho{Relatório Técnico}
%    {Relatório Técnico}
%    {Dissertação}
%    {Tese}
%    {Monografia}

% Marcar Sim para as partes que irão compor o documento pdf
%-----------------------------------------------------------------------------
 \providecommand{\terCapa}{Sim}
 \providecommand{\terFolhaRosto}{Nao}
 \providecommand{\terTermoAprovacao}{Nao}
 \providecommand{\terDedicatoria}{Nao}
 \providecommand{\terFichaCatalografica}{Nao}
 \providecommand{\terEpigrafe}{Nao}
 \providecommand{\terAgradecimentos}{Nao}
 \providecommand{\terErrata}{Nao}
 \providecommand{\terListaFiguras}{Nao}
 \providecommand{\terListaTabelas}{Nao}
 \providecommand{\terSiglasAbrev}{Nao}
 \providecommand{\terResumos}{Sim}
 \providecommand{\terSumario}{Sim}
 \providecommand{\terAnexo}{Nao}
 \providecommand{\terApendice}{Nao}
 \providecommand{\terIndiceR}{Nao}
%-----------------------------------------------------------------------------

%\titulo{Desenvolvimentos em Processamento de Dados Gravimetricos com Aplicação ao Estudo do Orógeno Dom Feliciano e da Porção Sul do Orógeno Ribeira}
\titulo{Relatório -- Prática Cálculo Do Geoide}
\autor{Eros Kerouak Cordeiro Pereira}
\local{Curitiba}
\data{2025} %Apenas ano 4 dígitos

% Orientador ou Orientadora
\orientador{}
%Prof Emílio Eiji Kavamura, MSc}
\orientadora{}
% Pode haver apenas uma orientadora ou um orientador
% Se houver os dois prevalece o feminino.

% Em termos de coorientação, podem haver até quatro neste modelo
% Sendo 2 mulhere e 2 homens.
% Coorientador ou Coorientadora
\coorientador{}%Prof Morgan Freeman, DSc}
\coorientadora{}

% Segundo Coorientador ou Segunda Coorientadora
\scoorientador{}
%Prof Jack Nicholson, DEng}
\scoorientadora{}
%Prof\textordfeminine~Ingrid Bergman, DEng}
% ----------------------------------------------------------
\addbibresource{referencias.bib}

% ----------------------------------------------------------
\instituicao{Universidade Federal do Paraná}

\def \ImprimirSetor{Setor de Ciências da Terra}%
%Setor de Tecnologia}

\def \ImprimirProgramaPos{MÉTODOS FÍSICOS EM GEODÉSIA (CGEO7023)}%Programa de Pós Graduação em Engenharia de Construção Civil}

\def \ImprimirCurso{}%
%Curso de Engenharia Civil}

\preambulo{
Este trabalho foi desenvolvido como parte das atividades da disciplina \textbf{MÉTODOS FÍSICOS EM GEODÉSIA (CGEO7023)}, do Programa de Pós-Graduação em Ciências Geodésicas da Universidade Federal do Paraná. O objetivo central da atividade foi aplicar, em ambiente computacional, os conceitos teóricos e metodológicos discutidos em aula, com ênfase em técnicas de remoção, interpolação e integração de anomalias de gravidade para a determinação de modelos residuais de geoide e quase-geoide.

A realização deste estudo permitiu não apenas a implementação prática de algoritmos de Colocação por Mínimos Quadrados, função de Stokes modificada e efeito indireto da topografia, mas também a análise crítica do desempenho dos modelos frente a marcos de rede de nivelamento. Dessa forma, a atividade cumpriu papel duplo: consolidar a compreensão conceitual dos métodos e fornecer experiência direta com fluxos reprodutíveis de processamento, aproximando a teoria à prática profissional e científica na área da geodésia física.}

%do grau de Bacharel em Expressão Gráfica no curso de Expressão Gráfica, Setor de Exatas da Universidade Federal do Paraná}

%-----------------------------------------------------------------------------

\newcommand{\imprimirCurso}{}
%Programa de P\'os Gradua\c{c}\~ao em Engenharia da Constru\c{c}\~ao Civil}

\newcommand{\imprimirDataDefesa}{
09 de Dezembro de 2018}

\newcommand{\imprimircdu}{
02:141:005.7}

% ----------------------------------------------------------
\newcommand{\imprimirerrata}{
Elemento opcional da \cites[4.2.1.2]{NBR14724:2011}. Exemplo:

\vspace{\onelineskip}

FERRIGNO, C. R. A. \textbf{Tratamento de neoplasias ósseas apendiculares com
reimplantação de enxerto ósseo autólogo autoclavado associado ao plasma
rico em plaquetas}: estudo crítico na cirurgia de preservação de membro em
cães. 2011. 128 f. Tese (Livre-Docência) - Faculdade de Medicina Veterinária e
Zootecnia, Universidade de São Paulo, São Paulo, 2011.

\begin{table}[htb]
\center
\footnotesize
\begin{tabular}{|p{1.4cm}|p{1cm}|p{3cm}|p{3cm}|}
  \hline
   \textbf{Folha} & \textbf{Linha}  & \textbf{Onde se lê}  & \textbf{Leia-se}  \\
    \hline
    1 & 10 & auto-conclavo & autoconclavo\\
   \hline
\end{tabular}
\end{table}}

% Comandos de dados - Data da apresentação
\providecommand{\imprimirdataapresentacaoRotulo}{}
\providecommand{\imprimirdataapresentacao}{}
\newcommand{\dataapresentacao}[2][\dataapresentacaoname]{\renewcommand{\dataapresentacao}{#2}}

% Comandos de dados - Nome do Curso
\providecommand{\imprimirnomedocursoRotulo}{}
\providecommand{\imprimirnomedocurso}{}
\newcommand{\nomedocurso}[2][\nomedocursoname]
  {\renewcommand{\imprimirnomedocursoRotulo}{#1}
\renewcommand{\imprimirnomedocurso}{#2}}


% ----------------------------------------------------------
\newcommand{\AssinaAprovacao}{

\assinatura{%\textbf
   {Professora} \\ UFPR}
   \assinatura{%\textbf
   {Professora} \\ ENSEADE}
   \assinatura{%\textbf
   {Professora} \\ TIT}
   %\assinatura{%\textbf{Professor} \\ Convidado 4}

   \begin{center}
    \vspace*{0.5cm}
    %{\large\imprimirlocal}
    %\par
    %{\large\imprimirdata}
    \imprimirlocal, \imprimirDataDefesa.
    \vspace*{1cm}
  \end{center}
  }

% ----------------------------------------------------------
%\newcommand{\Errata}{%\color{blue}
%Elemento opcional da \textcite[4.2.1.2]{NBR14724:2011}. Exemplo:
%}

% ----------------------------------------------------------
\newcommand{\EpigrafeTexto}{%\color{blue}
\textit{"O presente estado do universo é o efeito do seu estado anterior, e a causa do que se seguirá." \\
		(Pierre-Simon Laplace)}
}

% ----------------------------------------------------------
\newcommand{\ResumoTexto}{%\color{blue}

Modelos regionais de geoide dependem da compatibilização entre comprimentos de onda longos, bem descritos por modelos globais do geopotencial, e feições curtas, resolvidas por dados terrestres ruidosos e afetados pela topografia. Este trabalho tem por objetivo construir e validar uma superfície residual de geoide estatística e fisicamente consistente com marcos de rede de nivelamento. A metodologia emprega (i) remoção da contribuição do modelo GOCO06s (até grau 300) e do efeito indireto de Helmert para definir anomalias residuais de Helmert; (ii) estimação da covariância empírica e ajuste isotrópico polinomial de grau 9 até $4^\circ$; (iii) interpolação por Colocação por Mínimos Quadrados (LSC) com modelagem de ruído de $1{,}0$\,mGal e blocagem espacial; (iv) integração por PVCG com função de Stokes modificada, compatível com o grau removido; e (v) cálculo do efeito indireto (EI) em pontos e em grade, seguido de validação externa frente a marcos GNSS–nivelamento.

Os principais resultados indicam aderência elevada do modelo de covariância ($R^{2}=0{,}9549$), estabilidade da LSC e geração de uma grade residual contínua para integração. Na validação do geoide, obtêm-se viés de $-0{,}434$\,m, desvio-padrão de $0{,}144$\,m e $\mathrm{RMSE}=0{,}456$\,m; a remoção do viés reduziria o $\mathrm{RMSE}$ ao nível da dispersão ($\sim 0{,}144$\,m). O EI mostra contribuição topográfica sistemática e constitui componente adicional a ser combinado ao campo residual na avaliação final.

Conclui-se que as hipóteses sobre (H1) covariância polinomial, (H2) truncagem espectral na Stokes modificada e (H3) nível de ruído de $1{,}0$\,mGal são suportadas pelos produtos e métricas obtidos, enquanto (H4)—redução de enviesamento pela inclusão explícita do EI—permanece parcialmente aberta e requer testes adicionais. O fluxo apresentado torna explícita a ligação entre problemas (heterogeneidade, ruído, topografia) e soluções (remoção espectral, LSC, Stokes modificada, EI).


}

\newcommand{\PalavraschaveTexto}{%\color{blue}
}

% ----------------------------------------------------------
\newcommand{\AbstractTexto}{%\color{blue}

This work presents a methodology for integrating terrestrial and satellite gravity data for geophysical modeling of the subsurface, with application to the Southern Mantiqueira Province in southern Brazil. The approach adopts the \textit{Remove--Compute--Restore} (RCR) method, traditionally used for geoid determination, here adapted for geophysical purposes to enable the combination of multiscale information. Terrestrial gravity stations from various institutional sources were compiled and standardized, complemented by satellite data from the GOCO06s model, provided both as a regular grid and at the terrestrial observation points. The differences between terrestrial and satellite data were interpolated using the equivalent sources method, with depth and damping parameters calibrated through \textit{k-fold} cross-validation. The integrated grid was obtained by adding the interpolated differences to the reference satellite model. Spectral treatment involved the use of smoothed Digital Elevation Models (DEMs) for satellite data and original DEMs for terrestrial data, with prism-based discretization of the topography and sedimentary layers for the application of terrain and \textit{Sediment Stripping} corrections. These corrections were applied directly to the gravity disturbance, yielding the Bouguer Disturbance, which provides a more refined interpretation of the gravity field, less affected by geoid undulations and more sensitive to deep density variations. Finally, the Moho crustal interface was estimated based on the Airy hypothesis and modeled using tesseroids for the implementation of the isostatic correction.

}
% ---
\newcommand{\KeywordsTexto}{%\color{blue}
Equivalent sources. Sediment Stripping. Southern Mantiqueira Province.
}

% ----------------------------------------------------------
\newcommand{\Resume}
{%\color{blue}
%Il s'agit d'un résumé en français.
}
% ---
\newcommand{\Motscles}
{%\color{blue}
 %latex. abntex. publication de textes.
}

% ----------------------------------------------------------
\newcommand{\Resumen}
{%\color{blue}
%Este es el resumen en español.
}
% ---
\newcommand{\Palabrasclave}
{%\color{blue}
%latex. abntex. publicación de textos.
}

% ----------------------------------------------------------
\newcommand{\AgradecimentosTexto}{%\color{blue}
Os agradecimentos principais são direcionados à Gerald Weber, Miguel Frasson, Leslie H. Watter, Bruno Parente Lima, Flávio de  Vasconcellos Corrêa, Otavio Real Salvador, Renato Machnievscz\footnote{Os nomes dos integrantes do primeiro
projeto abn\TeX\ foram extraídos de \url{http://codigolivre.org.br/projects/abntex/}} e todos aqueles que contribuíram para que a produção de trabalhos acadêmicos conforme as normas ABNT com \LaTeX\ fosse possível.

Agradecimentos especiais são direcionados ao Centro de Pesquisa em Arquitetura da Informação\footnote{\url{http://www.cpai.unb.br/}} da Universidade de Brasília (CPAI), ao grupo de usuários
\emph{latex-br}\footnote{\url{http://groups.google.com/group/latex-br}} e aos novos voluntários do grupo \emph{\abnTeX}\footnote{\url{http://groups.google.com/group/abntex2} e
\url{http://abntex2.googlecode.com/}}~que contribuíram e que ainda
contribuirão para a evolução do \abnTeX.

Os agradecimentos principais são direcionados à Gerald Weber, Miguel Frasson, Leslie H. Watter, Bruno Parente Lima, Flávio de Vasconcellos Corrêa, Otavio Real Salvador, Renato Machnievscz\footnote{Os nomes dos integrantes do primeiro
projeto abn\TeX\ foram extraídos de \url{http://codigolivre.org.br/projects/abntex/}} e todos aqueles que contribuíram para que a produção de trabalhos acadêmicos conforme as normas ABNT com \LaTeX\ fosse possível.
}

% ----------------------------------------------------------
\newcommand{\DedicatoriaTexto}{%\color{blue}
\textit{ Este trabalho é dedicado às crianças adultas que,\\
   quando pequenas, sonharam em se tornar cientistas.}
	}
